\documentclass[12pt]{article}
\usepackage[utf8]{inputenc}
\usepackage[T1]{fontenc}
\usepackage[margin=1in]{geometry}
\usepackage{setspace}
\usepackage{xcolor}
\usepackage[colorlinks=true, urlcolor=blue!60!black, linkcolor=blue!60!black]{hyperref}
\setlength{\parindent}{2em}
\setlength{\parskip}{1.25em}
\renewcommand{\baselinestretch}{1.0}
\begin{document}
\begin{center}
{\LARGE \textit{Florida Prepaid Postsecondary Education Expense Board v.\ College Savings Bank}, 527 U.S.\ 627 (1999)\par}
\vspace{0.8em}
{\large Case Brief --- Four-Page Edition\par}
\end{center}
\vspace{1em}

\section{Memory Jogger}

Applying \textit{City of Boerne}'s ``congruence and proportionality'' test, the Court struck down the Patent Remedy Act's abrogation of state sovereign immunity under \S~5 of the Fourteenth Amendment, because Congress compiled virtually no record of unremedied, nonnegligent state patent infringement (\href{https://www.courtlistener.com/opinion/118316/florida-prepaid-postsecondary-education-expense-board-v-college-savings/}{CourtListener}).

\section{Detailed Case Facts}

College Savings Bank, a New Jersey savings bank, sold ``CollegeSure CDs,'' annuity-like certificates designed to cover future college tuition, and held a patent on its financing methodology. The Florida Prepaid Postsecondary Education Expense Board, an arm of the State of Florida, administered competing tuition prepayment contracts; College Savings alleged that it directly, contributorily, and willfully infringed the patent. After the Federal Circuit held in \textit{Chew v.\ California} (1990) that the Patent Act lacked the clear statement needed to abrogate state immunity, Congress passed the Patent Remedy Act of 1992, defining ``whoever'' to include States and declaring that States ``shall not be immune'' from federal infringement suits. Congress invoked the Patent Clause, the Commerce Clause, and \S~5 of the Fourteenth Amendment.

\section{Procedural History}

College Savings sued in the District of New Jersey in 1994 (a companion Lanham Act case was decided the same day). After \textit{Seminole Tribe} (1996) held that Congress cannot abrogate immunity under Article~I, Florida Prepaid moved to dismiss; the United States intervened to defend the Act. The district court denied the motion, and the Federal Circuit affirmed, holding patents to be Fourteenth Amendment property and the Act a proportionate \S~5 measure. The Supreme Court reversed and remanded.

\section{Judicial Votes}

Reversed, \textbf{5--4} (decided June 23, 1999).
\begin{itemize}
\item \textbf{Majority:} Chief Justice Rehnquist, joined by Justices O'Connor, Scalia, Kennedy, and Thomas.
\item \textbf{Dissent:} Justice Stevens, joined by Justices Souter, Ginsburg, and Breyer.
\end{itemize}

\section{Holding}

Congress clearly expressed its intent to abrogate, and patents are property under the Due Process Clause, but Article~I cannot support abrogation and the Act is not ``appropriate'' \S~5 legislation. Valid \S~5 legislation must respond to a pattern of conduct violating the Fourteenth Amendment (here, intentional or reckless infringement coupled with absent or inadequate state remedies) and show ``congruence and proportionality'' between injury and means. The Court's list of narrowing features Congress \textit{could} have adopted is dicta, and it did not reach the Just Compensation Clause.

\section{Analysis of Opinions}

\textbf{Majority (Rehnquist).} With Article~I foreclosed by \textit{Seminole Tribe}, only \S~5 remained, and under \textit{City of Boerne} it is a remedial power to enforce, not redefine, the Amendment. First, the record showed virtually no pattern of state infringement: the House Report cited two examples, and only eight suits against States had been filed in 110 years. Second, infringement alone does not violate due process: under \textit{Parratt} and \textit{Zinermon}, a deprivation occurs only when the State offers no adequate remedy, and Congress barely examined state remedies, focusing instead on uniformity, an Article~I concern. Under \textit{Daniels v.\ Williams}, negligence is no deprivation, yet infringement is strict liability and most state infringement appeared innocent. Third, the Act was disproportionate: it exposed all 50 States to suit for every form of infringement, indefinitely, without limiting itself to willful infringement or States lacking remedies.

\textbf{Dissent (Stevens).} The Act was ``a paradigm of an appropriate exercise of Congress' \S~5 power.'' Because federal courts have exclusive patent jurisdiction, Congress reasonably assumed state remedies were nonexistent or inadequate. The majority dismissed the Act on its face though College Savings alleged \textit{willful} infringement, and it was ``quite unfair'' to fault the 1992 Congress for not building a record for a test announced five years later. Unlike RFRA, the Act altered no substantive state law, and its reach scaled with actual state infringement. Stevens also renewed his dissent from \textit{Seminole Tribe}.

\section{Examples: Future Applications}

\begin{enumerate}
\item \textbf{Same-Side: Copyright Abrogation.} The Copyright Remedy Clarification Act rested on a thin anecdotal record and covered innocent infringement; it fails congruence and proportionality, as in \textit{Allen v.\ Cooper} (2020).
\item \textbf{Same-Side: ADEA Damages.} Subjecting States to age-discrimination damages without a record of irrational state discrimination is invalid, as in \textit{Kimel} (2000).
\item \textbf{Opposite-Side: Family Leave.} The FMLA's family-care provision, backed by an extensive record of gender stereotyping and narrowly tailored, is valid, as in \textit{Hibbs} (2003).
\item \textbf{Opposite-Side: A Narrow Patent Statute.} An Act reaching only willful infringement under state policy where no adequate state remedy exists would likely be upheld under the Court's own roadmap.
\item \textbf{Fence-Sitter: Conditional Spending.} Conditioning federal research grants on a State's waiver of immunity avoids abrogation, but its validity under \textit{South Dakota v.\ Dole} turns on clarity, germaneness, and coercion.
\end{enumerate}

\section{Critique}

\textbf{Federalism and text.} Defenders praise the decision for stopping \S~5 end-runs around \textit{Seminole Tribe}. Textualists such as John F.\ Manning note that the Eleventh Amendment's text bars only diversity suits, so the doctrine rests on nontextual ``presuppositions''; others criticize the facial invalidation despite alleged willfulness and the quasi-administrative record requirement.

\textbf{Progressive critique.} Erwin Chemerinsky, Robert Post, and Reva Siegel see a ``juricentric Constitution'' that strips Congress of \S~5 discretion, and commentators note the asymmetry: States enforce their own patents in federal court but escape liability when they infringe.

\section{Key Quotations}

\begin{quote}
``We thus held that for Congress to invoke \S~5, it must identify conduct transgressing the Fourteenth Amendment's substantive provisions, and must tailor its legislative scheme to remedying or preventing such conduct.''
\end{quote}
\noindent --- Chief Justice Rehnquist, for the Court, 527 U.S.\ at 639.

\begin{quote}
``Thus, under the plain terms of the Clause and the clear import of our precedent, a State's infringement of a patent, though interfering with a patent owner's right to exclude others, does not by itself violate the Constitution. Instead, only where the State provides no remedy, or only inadequate remedies, to injured patent owners for its infringement of their patent could a deprivation of property without due process result.''
\end{quote}
\noindent --- Chief Justice Rehnquist, for the Court, 527 U.S.\ at 643.

\begin{quote}
``The need for uniformity in the construction of patent law is undoubtedly important, but that is a factor which belongs to the Article I patent-power calculus, rather than to any determination of whether a state plea of sovereign immunity deprives a patentee of property without due process of law.''
\end{quote}
\noindent --- Chief Justice Rehnquist, for the Court, 527 U.S.\ at 645.

\end{document}
